\section{Conclusions}\label{sec:conclusion}

This study constructed a four-dimension DFA index covering digital infrastructure, digital industry, digital finance, and digital human capital, and a three-dimension RCD index covering economic, social, and urban-rural coordination, for 30 Chinese provinces from 2013 to 2021. Using a Hansen panel threshold model, it tested whether the effect of DFA on RCD depends on economic development and financial development, and it examined which dimension carries the association and whether the estimates survive spatial dependence.

DFA is positively associated with RCD. The baseline two-way fixed-effects estimate is 0.241 and the association holds across alternative weighting, sample restrictions, and winsorization. It is economically meaningful: a one-standard-deviation increase in DFA is associated with an increase of about 0.043 points in RCD, or 0.29 standard deviations.

The association rises with economic development. Above the estimated threshold of 10.589 in log GDP per capita (about CNY 39,700), the marginal association is 0.677 against 0.506 below, consistent with absorptive capacity. A financial-development threshold indicates declining returns at very high levels of financial depth, driven by a small group of financially deep municipalities. Both indices are spatially clustered; allowing for spatial dependence leaves the threshold location unchanged, so we read it as a development gradient rather than a sharp break.

Two findings extend the usual account. Within the composite, digital finance is the only dimension with an independent association, and it also shows the most stable positive slope across regimes. Industrial digitalization and digital human capital, by contrast, switch from negative below the threshold to positive above it, which is what an absorptive-capacity argument predicts. The regional analysis shows the association concentrated in the west, where digital infrastructure and financial access start low.

These findings imply a differentiated policy. Digital factor allocation should be sequenced with economic development: basic connectivity and digital literacy first in low-income provinces, deeper integration of digital technologies into production and public services in higher-income ones. The west should be a priority for digital infrastructure and digital finance investment, paired with human capital development and market reform. Because the short-run association with the urban-rural income ratio is positive, complementary rural digital services and targeted skills training are needed to share the gains. Advanced regions should focus on the quality and direction of digital investment rather than its quantity.

Several limitations remain. The analysis uses provincial data and does not capture within-province heterogeneity; the indices depend on indicator and weighting choices; and the estimates are conditional associations rather than identified causal effects. The threshold is sensitive to the treatment of trends, year effects, and spatial dependence. Future work can address these points with city-level data, external instruments, dynamic spatial threshold models, and a DFA index extended to data governance, digital security, and the quality of digital public services.

Digitalization is not a universal remedy for regional disparities. Its payoff depends on the stage of economic development, the depth of the financial system, and regional absorptive capacity, and a differentiated approach that combines digital investment with complementary economic and social reforms is more likely to deliver inclusive coordination gains.
